\chapter{Résultats}\label{ch:s1:earnings}

Une entreprise publie des résultats nettement supérieurs à ce que le marché attendait, son action bondit, puis, pendant des semaines, le cours continue de dériver dans le même sens. La
\omterm{def:s1:earnings:pead}{dérive post-annonce des résultats} (PEAD) est l'une des plus anciennes anomalies de la recherche comptable et l'une des plus étudiées~; elle est aussi celle dont le bilan public
s'est retourné contre elle. Chordia et ses coauteurs l'ont trouvée concentrée dans les actions les moins liquides, où les coûts de transaction absorbaient 70 à 100~\% de ses profits sur le papier,
et Martineau la juge inexistante dans les grandes actions américaines depuis 2006. Le marché synthétique en fournit une version épurée~: 1.2~\% de dérive par écart type de
surprise, étalée sur soixante jours. Négociée à partir du lendemain de l'annonce et conservée pendant soixante jours, elle donne un ratio de Sharpe de 2.66 après coûts~; conservée pendant
cinq jours, $-0.77$~; et lorsque la dérive perd 70~\% de son amplitude, comme dans le marché synthétique du livre~7, chapitre~13, le ratio de Sharpe du portefeuille à soixante jours sur les
quatre dernières années passe de 3.24 à 0.75. Ce chapitre traite des transactions autour des résultats. La construction s'appelle \texttt{firm.earnstrat}.

\section{Dérive post-annonce des résultats}

\begin{definition}[Étude d'événement, fenêtre d'événement]\label{def:s1:earnings:event}
Une \emph{étude d'événement}\index{etude@étude d'événement} (event study) mesure le rendement anormal moyen des titres en \emph{temps d'événement}, par rapport à la date d'un événement
d'un type donné, sur une \emph{fenêtre d'événement}\index{fenêtre d'événement} (event window) de jours avant et après celui-ci, afin de séparer l'effet de l'événement de celui du marché.
\end{definition}

\begin{definition}[Dérive post-annonce des résultats]\label{def:s1:earnings:pead}
La \emph{dérive post-annonce des résultats}\index{dérive post-annonce des résultats} (post-earnings-announcement drift) est la tendance des rendements anormaux d'un titre à se poursuivre dans le sens de
sa surprise de résultats, pendant des semaines après l'annonce.
\end{definition}

Ball et Brown (1968) et Bernard et Thomas (1989) sont les études classiques~; Chan, Jegadeesh et Lakonishok ont établi que les surprises de résultats passées et les rendements passés
prédisent chacun de fortes dérives des rendements futurs, une fois neutralisé l'effet de l'autre, et que les prévisions des analystes réagissent elles aussi lentement aux nouvelles~: un marché
qui absorbe l'information graduellement. Le marché synthétique (livre~7, chapitre~5) reproduit exactement cela~: à chaque annonce, une surprise normale centrée réduite, un saut de
trois volatilités spécifiques quotidiennes par unité de surprise, et une dérive de 1.2~\% par unité sur les soixante jours suivants. Les entreprises publient quatre fois par an (1.59~\% des
jours-titres sont des jours d'annonce). La \Cref{fig:s1:earnings:car} est l'\omterm{def:s1:earnings:event}{étude d'événement} du chapitre~: pour des surprises supérieures à un écart type, le
rendement corrigé du marché est de 6.94~\% le jour de l'annonce et de 7.75~\% après soixante jours~; en deçà de moins un, $-6.93~\%$ et $-8.56~\%$.

\begin{omfigure}
\begin{tikzpicture}
\begin{axis}[width=11cm,height=5.2cm,xlabel={\small jours depuis l'annonce},ylabel={\small rendement anormal cumulé (\%)},
  tick label style={font=\footnotesize},xmin=-5,xmax=60,ymin=-10,ymax=10,grid=major,grid style={gray!25},
  legend style={font=\scriptsize,at={(0.97,0.5)},anchor=east},table/col sep=comma]
\addplot[omDef,very thick] table[x=day,y=positive_pct]{figdata/strategies-1/07-earnings/car.csv};
\addplot[omThm,very thick] table[x=day,y=negative_pct]{figdata/strategies-1/07-earnings/car.csv};
\addplot[black!50,dashed] coordinates {(-5,0) (60,0)};
\legend{surprise supérieure à $+1$ é.-t.,surprise inférieure à $-1$ é.-t.}
\end{axis}
\end{tikzpicture}

\omcaption{\omterm{def:s1:earnings:event}{Étude d'événement} des annonces du marché synthétique~: rendement corrigé du marché cumulé moyen, de cinq jours avant à soixante jours après, pour
grandes surprises positives et grandes surprises négatives. Données~: \texttt{s1\_earnings.car}.}\label{fig:s1:earnings:car}
\end{omfigure}

L'\omterm{def:s1:earnings:event}{étude d'événement} montre aussi le lendemain~: le groupe positif rend 0.33 point, passant de 6.94~\% à 6.61~\%. C'est le retournement d'un jour inscrit dans le marché synthétique, celui
du chapitre~2 qui agit sur le saut de l'annonce, et il décide du premier choix de mise en œuvre ci-dessous.

\section{Révisions des analystes}

\begin{definition}[Stratégie de révisions des résultats]\label{def:s1:earnings:revision}
Une \emph{stratégie de révisions des résultats}\index{stratégie de révisions des résultats} (earnings-revision strategy) achète les titres dont les analystes ont récemment relevé leurs prévisions de résultats et vend
ceux dont les prévisions ont été abaissées, au motif que les prévisions, et les cours avec elles, s'ajustent progressivement aux nouvelles.
\end{definition}

La même lenteur qui produit la dérive se retrouve dans les prévisions~: Chan, Jegadeesh et Lakonishok ont constaté que les prévisions des analystes réagissent lentement aux nouvelles passées,
surtout pour les moins bonnes performances. Un signal de révision est la variation de la prévision consensuelle sur un mois, normalisée par le cours ou par la dispersion des
prévisions (livre~7, chapitre~11). Il faut disposer d'un historique de prévisions individuelles des analystes, daté au point près, ce qui relève d'un jeu de données commercial~; le marché synthétique
ne compte aucun analyste, si bien que la fiche de stratégie sur les révisions du chapitre repose sur le bilan publié.

\section{Orientations et pré-annonces}

Les nouvelles sur les résultats n'arrivent pas seulement à la date prévue. Les entreprises donnent des orientations sur leurs résultats futurs, les mettent à jour entre deux publications, et parfois
pré-annoncent un résultat qui manquera ou dépassera les attentes, généralement dans les semaines qui précèdent la publication trimestrielle. Chacun de ces cas est un événement doté de sa propre fenêtre,
et la question de la dérive s'applique à chacun~: le cours se déplace-t-il entièrement le jour même, ou continue-t-il~? La mise en œuvre est la même que pour les annonces (un calendrier
des événements connus à la date près et un portefeuille en \omterm{def:s1:earnings:event}{temps d'événement}), et le problème des données est plus délicat~: les orientations sont publiées dans des communiqués de presse et lors des conférences téléphoniques, et non dans un
dépôt normalisé, et leur horodatage doit correspondre au moment où elles sont devenues publiques. La fiche de stratégie sur les pré-annonces indique qu'aucun bilan de performance n'a été
vérifié pour ce livre.

\section{L'angle des options le jour de l'annonce}

\begin{definition}[Prime d'annonce]\label{def:s1:earnings:premium}
La \emph{prime d'annonce}\index{prime d'annonce} (announcement premium) est le rendement moyen plus élevé des titres dans les jours entourant leurs
annonces de résultats programmées, par rapport aux mêmes titres à d'autres moments.
\end{definition}

Savor et Wilson ont établi que les entreprises dont la publication était programmée affichaient un rendement anormal annualisé de 9.9~\%, et l'ont expliqué par une prime de risque~: les annonces
renseignent aussi les investisseurs sur les autres entreprises et sur le marché, si bien que la covariance entre les nouvelles d'une entreprise et celles du marché s'envole autour d'elles. Le marché synthétique
n'intègre aucune \omterm{def:s1:earnings:premium}{prime d'annonce}, et le portefeuille des titres annoncés chaque jour, acheté à la clôture de la veille, rapporte 1.8 points de base par jour au-dessus du marché avec une
statistique $t$ de 0.59~: rien.

Ce que le marché synthétique possède en revanche, c'est le saut de l'annonce~: le rendement corrigé en valeur absolue le jour d'une annonce vaut 3.02 fois celui
d'un jour ordinaire. Une option qui expire juste après une annonce doit intégrer ce mouvement supplémentaire, si bien que sa volatilité implicite devrait être plus élevée que celle d'une
option qui expire juste avant, et devrait baisser une fois l'annonce publiée. Un straddle acheté avant l'annonce est un pari sur un mouvement qui dépassera ce que les options ont intégré, et un straddle vendu est
le pari inverse. Aucun des deux n'est un gain sans risque~: ce sont tous deux des transactions sur la taille du saut relativement à son prix (livre~9, chapitre~4, transactions sur la volatilité réalisée face à la volatilité implicite autour des événements).

\section{Mise en œuvre autour des dates d'événement}

\begin{center}
\small
\begin{tabular}{@{}p{8.4cm}ccc@{}}
\toprule
portefeuille de dérive, marché synthétique, années 3 à 10 & 5 jours & 20 jours & 60 jours\\
\midrule
acheté à la clôture de l'annonce~: ratio de Sharpe après coûts & $-3.72$ & 0.02 & 2.03\\
acheté le lendemain~: ratio de Sharpe avant, après coûts & 1.38, $-0.77$ & 2.64, 1.51 & 3.31, 2.66\\
acheté le lendemain~: rendement net par an~; rotation par an & $-4.7~\%$~; 65.2 & 4.6~\%~; 17.2 & 5.5~\%~; 6.6\\
sur le rendement du jour d'annonce plutôt que sur la surprise & $-0.85$ & 0.85 & 2.03\\
dérive réduite de 70~\% à partir de l'année 7~: ratio de Sharpe sur les années 7 à 10 & $-1.19$ & $-0.51$ & 0.75\\
\bottomrule
\end{tabular}
\end{center}

Trois choix déterminent le résultat. \emph{Quand entrer}~: acheter à la clôture de l'annonce revient à garder la position pendant le renversement du lendemain du saut, ce qui
rend négatif le ratio de Sharpe brut du portefeuille à cinq jours~; attendre un jour coûte un jour de dérive et évite ce problème. \emph{Combien de temps conserver}~: la dérive s'accumule lentement
(1.2~\% par unité de surprise sur soixante jours, soit environ deux points de base par jour), si bien qu'une courte période de détention paie un aller-retour complet pour quelques jours de cette dérive~; le
portefeuille à cinq jours échange 65 fois sa taille brute par an, et ses coûts, 13~\% par an, dépassent son rendement brut. \emph{Ce qu'il faut mesurer}~: la surprise publiée l'emporte sur le
rendement du jour d'annonce comme signal (2.66 contre 2.03 à soixante jours), car ce rendement mêle la surprise au bruit spécifique du jour~; les portefeuilles réels
combinent les deux et y ajoutent les révisions.

Et la dérive peut se réduire. Sur le marché synthétique dont la dérive perd 70~\% de sa taille à partir de l'année 7 (\texttt{MarketConfig.pead\_break}), le livre à soixante jours
gagne 4.8~\% par an avant la rupture et 1.4~\% après, et son ratio de Sharpe passe de 2.18 à 0.75~; le livre à vingt jours passe de 1.48 à $-0.51$. Le
constat public va dans le même sens~: Chordia et coauteurs, 0.04~\% par mois dans les actions les plus liquides, et la dérive de Martineau qui a disparu des grandes actions après
2006.

\begin{omfigure}
\begin{tikzpicture}
\begin{axis}[ybar,bar width=10pt,width=10cm,height=5.2cm,xlabel={\small période de détention (jours)},ylabel={\small ratio de Sharpe après coûts},
  tick label style={font=\footnotesize},symbolic x coords={5,20,60},xtick=data,ymin=-4.5,ymax=4,grid=major,grid style={gray!25},
  legend style={font=\scriptsize,at={(0.5,1.03)},anchor=south,legend columns=1},table/col sep=comma,enlarge x limits=0.25]
\addplot[fill=omThm!60,draw=omThm] table[x=hold,y=same_day]{figdata/strategies-1/07-earnings/holds.csv};
\addplot[fill=omDef!70,draw=omDef] table[x=hold,y=next_day]{figdata/strategies-1/07-earnings/holds.csv};
\addplot[fill=black!25,draw=black!60] table[x=hold,y=cut_late]{figdata/strategies-1/07-earnings/holds.csv};
\legend{acheté à l'annonce,acheté un jour plus tard,un jour plus tard~; dérive réduite (années 7--10)}
\end{axis}
\end{tikzpicture}

\omcaption{Ratio de Sharpe du livre sur la dérive après 10 points de base par unité échangée, selon la durée de détention et le jour d'entrée, et après que la dérive a perdu 70~\% de sa taille.
Données~: \texttt{s1\_earnings.run}.}\label{fig:s1:earnings:holds}
\end{omfigure}

\section{Fiches de stratégie}

\begin{strategyfile}{Dérive SUE}\label{strat:s1:earnings:sue}
\sfield{Qui paie, et pourquoi} Investisseurs qui sous-réagissent aux annonces de résultats et s'ajustent sur plusieurs semaines.
\sfield{Instruments et plateformes} Actions, en position longue sur les surprises positives et courte sur les négatives.
\sfield{Signal} Résultat inattendu standardisé (livre~7, chapitre~11) ou surprise par rapport au consensus, connue à l'annonce.
\sfield{Dimensionnement et exécution} Livre en \omterm{def:s1:earnings:event}{temps d'événement}, entré le lendemain, conservé jusqu'à soixante jours.
\sfield{Coûts} Le point faible de la stratégie~: la dérive se trouve dans des actions peu liquides, où les coûts sont les plus élevés.
\sfield{Comment elle s'éteint} Une formation des prix plus rapide le jour de l'annonce~; les coûts~; l'encombrement (crowding).
\sfield{Horizon, capacité, infrastructure} Quelques semaines~; faible capacité dans les actions où la dérive subsiste~; un calendrier des annonces et des données de surprises de résultats.
\sfield{Backtest en toute honnêteté} Horodatages des annonces (avant ou après la séance)~; consensus tel qu'il était~; coûts selon la liquidité.
\sfield{Sources} Chan, Jegadeesh et Lakonishok (1996)~; Chordia et al.\ (2009)~: 0.04~\% par mois dans les actions les plus liquides, 2.43~\% dans les moins liquides, les coûts
absorbent 70--100~\% des profits théoriques~; Martineau (2022)~: aucune dérive dans les grandes actions depuis 2006.
\end{strategyfile}

\begin{strategyfile}{Momentum des révisions}\label{strat:s1:earnings:revision}
\sfield{Qui paie, et pourquoi} Comme pour la dérive~: les analystes et les cours s'ajustent progressivement.
\sfield{Instruments et plateformes} Actions suivies par des analystes.
\sfield{Signal} La variation sur un mois des prévisions consensuelles, normalisée par le cours ou par la dispersion.
\sfield{Dimensionnement et exécution} Tris mensuels, ou signal dans un livre multifactoriel.
\sfield{Coûts} Rotation modérée.
\sfield{Comment elle s'éteint} Comme pour la dérive~; effet de troupeau dans les prévisions.
\sfield{Horizon, capacité, infrastructure} Mois~; un historique de prévisions point-in-time.
\sfield{Backtest en toute honnêteté} Prévisions datées de leur émission, et non de leur saisie dans la base~; prévisions périmées retirées.
\sfield{Sources} Chan, Jegadeesh et Lakonishok (1996) sur les prévisions lentes~; aucun chiffre de performance vérifié pour le présent ouvrage.
\end{strategyfile}

\begin{strategyfile}{Prime d'annonce}\label{strat:s1:earnings:premium}
\sfield{Qui paie, et pourquoi} Investisseurs qui exigent une prime pour le risque de marché porté par les entreprises qui font une annonce.
\sfield{Instruments et plateformes} Actions dont les annonces sont programmées.
\sfield{Signal} La date d'annonce, connue à l'avance.
\sfield{Dimensionnement et exécution} Position longue sur les actions qui annoncent dans les prochains jours, couverte par le marché.
\sfield{Coûts} Un aller-retour autour de chaque annonce.
\sfield{Comment elle s'éteint} Comme la prime est une prime de risque, elle se paie avec le risque~: des pertes lorsque les annonces apportent de mauvaises nouvelles pour le marché.
\sfield{Horizon, capacité, infrastructure} Jours~; un calendrier prospectif fiable.
\sfield{Backtest en toute honnêteté} Dates programmées telles qu'elles étaient connues à l'avance, et non les dates réalisées.
\sfield{Sources} Savor et Wilson (2016)~: un rendement anormal annualisé de 9.9~\% pour les sociétés dont la publication est programmée.
\end{strategyfile}

\begin{strategyfile}{Straddle avant les annonces}\label{strat:s1:earnings:straddle}
\sfield{Qui paie, et pourquoi} Celui qui a vendu les options trop bon marché, si le mouvement dépasse celui qui est intégré dans le prix~; l'acheteur paie sinon.
\sfield{Instruments et plateformes} Options sur actions cotées, arrivant à échéance juste après l'annonce.
\sfield{Signal} Le mouvement d'annonce implicite dans les options, comparé à une prévision du mouvement effectif.
\sfield{Dimensionnement et exécution} Straddles couverts en delta, achetés ou vendus avant l'annonce, clôturés après.
\sfield{Coûts} Fourchettes des options, élevées autour des événements.
\sfield{Comment elle s'éteint} Mouvements implicites qui évaluent correctement l'événement en moyenne.
\sfield{Horizon, capacité, infrastructure} Jours~; des données d'options et un modèle de mouvement.
\sfield{Backtest en toute honnêteté} Cours des options tels que cotés, cours acheteur et cours vendeur~; mouvement implicite de l'événement séparé de la volatilité ordinaire.
\sfield{Sources} Simulation de ce chapitre seulement (mouvements d'annonce 3.02 fois les mouvements ordinaires)~; aucun historique public de performance vérifié.
\end{strategyfile}

\begin{strategyfile}{Trading sur pré-annonces}\label{strat:s1:earnings:preannounce}
\sfield{Qui paie, et pourquoi} Sous-réaction aux changements de guidance et aux pré-annonces, si elle existe.
\sfield{Instruments et plateformes} Actions.
\sfield{Signal} Changements de guidance et pré-annonces, issus de communiqués de presse, horodatés à leur publication.
\sfield{Dimensionnement et exécution} Livre en \omterm{def:s1:earnings:event}{temps d'événement}, comme pour la dérive.
\sfield{Coûts} Comme pour la dérive.
\sfield{Comment elle s'éteint} Comme pour la dérive.
\sfield{Horizon, capacité, infrastructure} Semaines~; un flux d'information analysé pour en extraire la guidance.
\sfield{Backtest en toute honnêteté} Horodatages publics seulement~; aucune information avant la publication (chapitre~17, consacré aux nouvelles).
\sfield{Sources} Aucun historique de performance vérifié pour le présent ouvrage.
\end{strategyfile}

\section{Tutoriel~: soixante jours de dérive}\label{tut:s1:earnings:tut}

\begin{tutorial}
\textbf{Objectif.} Trader la dérive simulée en \omterm{def:s1:earnings:event}{temps d'événement}, choisir le jour d'entrée et la durée de détention, et observer la réduction de la dérive. \textbf{État final~:} le
tableau et la \cref{fig:s1:earnings:holds}.
\begin{tutsteps}
\item \textbf{Le livre en \omterm{def:s1:earnings:event}{temps d'événement}}~: valeurs ayant une forte surprise récente, à partir d'un délai après l'annonce, sur une fenêtre de détention.
\omcode{code/firm/earnstrat/firm_earnstrat.py}{34}{48}{Le livre en temps d'événement.}\label{lst:s1:earnings:book}
\item \textbf{L'\omterm{def:s1:earnings:event}{étude d'événement}}~: rendements anormaux cumulés en \omterm{def:s1:earnings:event}{temps d'événement}.
\omcode{code/firm/earnstrat/firm_earnstrat.py}{57}{67}{Rendements anormaux cumulés en temps d'événement.}\label{lst:s1:earnings:car}
\item \textbf{Lancer} \texttt{s1\_earnings.run(hold, signal, cut, delay)} pour 5, 20 et 60 jours, \texttt{car}, \texttt{moves}, \texttt{announcers} et
\texttt{fig\_earnings.py}.
\end{tutsteps}
\textbf{Pour aller plus loin.} Pondérer les positions selon l'ampleur de la surprise~; combiner la surprise et le rendement d'annonce~; restreindre le livre aux 500
valeurs les plus liquides et mesurer ce que vaut la dérive sur ce périmètre.
\end{tutorial}

\section{Construction~: stratégies sur les résultats}\label{bld:s1:earnings:earnstrat}

\begin{build}
\bfield{Rôle} Calendriers des événements, livres en \omterm{def:s1:earnings:event}{temps d'événement} et \omterm{def:s1:earnings:event}{études d'événement} pour les résultats et d'autres événements datés.
\bfield{Interface} \texttt{calendar(events, T, N)}, \texttt{event\_book(flag, surprise, listed, hold, threshold, delay)},
\texttt{announcement\_move(ret, flag)}, \texttt{event\_car(abret, flag, surprise, before, after, threshold)}.
\bfield{Règles} Événements datés à leur publication~; positions prises à la clôture qui suit le signal~; fenêtres de détention fixées à l'avance.
\bfield{Tests d'acceptation} \texttt{code/firm/earnstrat/tests/}~: le calendrier et le livre à la main, avec et sans délai~; le ratio des mouvements et la trajectoire
de l'événement sur un exemple simulé.
\bfield{Pour aller plus loin} SUE à partir du BPA point-in-time (livre~7, chapitre~11, \texttt{firm.fundpit})~; signaux de révisions~; le calendrier des annonces comme calendrier prospectif.
\end{build}

\begin{omsources}
\item R.~Ball et P.~Brown, «~An empirical evaluation of accounting income numbers~», \emph{Journal of Accounting Research} 6(2), 1968.
\item V.~L.~Bernard et J.~K.~Thomas, «~Post-earnings-announcement drift: delayed price response or risk premium?~», \emph{Journal of Accounting Research}
27, 1989.
\item L.~K.~C.~Chan, N.~Jegadeesh et J.~Lakonishok, «~Momentum strategies~», \emph{Journal of Finance} 51(5), 1996.
\item P.~Savor et M.~Wilson, «~Earnings announcements and systematic risk~», \emph{Journal of Finance} 71(1), 2016.
\item T.~Chordia, A.~Goyal, G.~Sadka, R.~Sadka et L.~Shivakumar, «~Liquidity and the post-earnings-announcement drift~», \emph{Financial Analysts
Journal} 65(4), 2009.
\item C.~Martineau, «~Rest in peace post-earnings announcement drift~», \emph{Critical Finance Review} 11(3--4), 2022.
\end{omsources}

\section{Exercices}

\begin{exercise}[$\star$]\label{exo:s1:earnings:1}
Avec une dérive de 1.2~\% par unité de surprise sur soixante jours, de combien dérive en tout une surprise de 1.5 écart-type, et quelle est cette dérive
par jour et par unité de surprise~?
\end{exercise}

\begin{exercise}[$\star$]\label{exo:s1:earnings:2}
Le livre à cinq jours échange 65.2 fois son brut par an, et le livre à soixante jours 6.6 fois. Que coûtent chacun, par an, dix points de base par unité échangée~?
\end{exercise}

\begin{exercise}[$\star$]\label{exo:s1:earnings:3}
Le marché synthétique inverse 5~\% du choc spécifique de chaque jour le lendemain. Que reprend cela sur un mouvement de 6.94~\% le jour de l'annonce~?
\end{exercise}

\begin{exercise}[$\star\star$]\label{exo:s1:earnings:4}
Pourquoi la surprise publiée est-elle ici un meilleur signal que le rendement du jour de l'annonce, et dans quels cas ce rendement pourrait-il être le meilleur~?
\end{exercise}

\begin{exercise}[$\star\star$]\label{exo:s1:earnings:5}
Pourquoi la date de l'annonce est-elle un signal légitime, alors que la surprise n'est connue qu'à l'annonce~? Que doit connaître un backtest de l'horodatage~?
\end{exercise}

\begin{exercise}[$\star\star$]\label{exo:s1:earnings:6}
Les mouvements le jour de l'annonce sont trois fois plus forts que les mouvements ordinaires. Qu'implique cela pour la volatilité implicite d'une option qui expire juste après l'annonce,
par rapport à une option qui expire juste avant~?
\end{exercise}

\begin{exercise}[$\star\star\star$]\label{exo:s1:earnings:7}
\emph{Programmation.} Exécuter \texttt{run(60, 'surprise', True)}. Donner les ratios de Sharpe avant et après la rupture, ainsi que les rendements annuels, et dire combien de temps
il faudrait observer pour détecter le changement (livre~7, chapitre~13).
\end{exercise}

\begin{exercise}[$\star\star\star$]\label{exo:s1:earnings:8}
\emph{Trouver la faille.} «~Notre backtest du SUE sur toutes les actions américaines depuis 1990 rapporte 2~\% par mois~; nous le ferons tourner sur le S\&P 500.~»
\end{exercise}

\section{Problème~: Soixante jours de dérive}

\begin{problem}[{Problème du week-end --- une anomalie exploitée, puis qui s'estompe}]\label{pb:s1:earnings:1}
La \omterm{def:s1:earnings:pead}{dérive post-annonce des résultats} introduite dans \texttt{firm.synthmkt}, et les travaux publiés.

\medskip\noindent\textbf{Partie~I --- L'effet.}
\begin{enumerate}
\item Définir une \omterm{def:s1:earnings:event}{étude d'événement}, une \omterm{def:s1:earnings:event}{fenêtre d'événement} et la dérive.
\item Qu'ont établi Chan, Jegadeesh et Lakonishok~?
\item Décrire la dérive introduite dans le marché synthétique et la fréquence de ses annonces.
\item Lire l'\omterm{def:s1:earnings:event}{étude d'événement}~: rendements le jour de l'annonce et sur soixante jours pour les deux groupes.
\end{enumerate}

\medskip\noindent\textbf{Partie~II --- Le livre.}
\begin{enumerate}[resume]
\item Décrire le livre en \omterm{def:s1:earnings:event}{temps d'événement}.
\item Donner les ratios de Sharpe à 5, 20 et 60 jours, achetés à l'annonce et un jour plus tard.
\item Pourquoi l'achat à la clôture du jour de l'annonce nuit-il~?
\item Pourquoi les détentions courtes perdent-elles après coûts~?
\end{enumerate}

\medskip\noindent\textbf{Partie~III --- Variantes.}
\begin{enumerate}[resume]
\item Surprise ou rendement du jour de l'annonce~: lequel l'emporte, et pourquoi~?
\item Qu'est-ce que la \omterm{def:s1:earnings:premium}{prime d'annonce}, et que montre le marché synthétique~?
\item Quelle est l'ampleur des mouvements à l'annonce, et qu'implique cela pour les options~?
\item Que sont les stratégies de révisions et celles qui précèdent l'annonce, et de quelles données ont-elles besoin~?
\end{enumerate}

\medskip\noindent\textbf{Partie~IV --- Le verdict.}
\begin{enumerate}[resume]
\item Énoncer le \emph{résultat nommé}~: le ratio de Sharpe du livre de dérive pour des détentions de 5, 20 et 60 jours, et sa chute après la rupture introduite.
\item Que disent les travaux publiés sur la dérive dans les actions liquides et de grande capitalisation~?
\item Où la chercher encore~?
\item Comment surveiller un livre de dérive en exploitation pour en détecter l'érosion~?
\item Quels sont les pièges liés à l'horodatage des données de résultats~?
\item Quelle fiche de stratégie faudrait-il lancer, et à quelle taille~?
\item Qu'enseigne la \omterm{def:s1:earnings:premium}{prime d'annonce} sur le risque et les anomalies~?
\item En une phrase~: qu'est-ce qu'une stratégie sur les résultats~?
\end{enumerate}
\end{problem}

\section{Questions d'entretien}

\begin{interviewq}[$\star$ \iqroles{chercheur}]\label{iq:s1:earnings:1}
Qu'est-ce que la \omterm{def:s1:earnings:pead}{dérive post-annonce des résultats} (PEAD), et pourquoi peut-elle exister~?
\end{interviewq}

\begin{interviewq}[$\star\star$ \iqroles{chercheur, développeur}]\label{iq:s1:earnings:2}
Comment construiriez-vous un signal de surprise de résultats point-in-time~? Que peut-il mal tourner~?
\end{interviewq}

\begin{interviewq}[$\star\star$ \iqroles{chercheur}]\label{iq:s1:earnings:3}
Une anomalie était forte dans les années 1990 et ne l'est plus aujourd'hui. Comment décidez-vous s'il faut continuer à l'exploiter~?
\end{interviewq}

\begin{interviewq}[$\star\star$ \iqroles{trader}]\label{iq:s1:earnings:4}
Comment négocieriez-vous une annonce de résultats avec des options~?
\end{interviewq}

\begin{interviewq}[$\star\star$ \iqroles{risque}]\label{iq:s1:earnings:5}
Quels risques porte un livre de résultats piloté par les événements et que n'a pas un livre factoriel~?
\end{interviewq}

\begin{interviewq}[$\star\star\star$ \iqroles{chercheur}]\label{iq:s1:earnings:6}
Une dérive de $d$ par unité de surprise s'accumule uniformément sur $H$ jours~; les coûts de transaction valent $c$ par unité échangée. Déterminez la durée de détention qui maximise le rendement net du livre
par unité de capital lorsque chaque position n'est ouverte et refermée qu'une fois.
\end{interviewq}
